\begin{table}[H]
\centering\centering\centering
\caption{Roster racial composition and win percentage, snap-weighted share, 2014-2025 \label{tab:roster-diversity-main}}
\centering
\resizebox{\ifdim\width>\linewidth\linewidth\else\width\fi}{!}{
\begin{tabular}[t]{lccccccccc}
\toprule
  & (1) & (2) & (3) & (4) & (5) & (6) & (7) & (8) & (9)\\
\midrule
Roster share Black (snap-weighted) & -0.684*** & -0.601** & -0.628** & -0.566** & -0.694*** &  & -0.692*** & -3.642 & \\
 & (0.209) & (0.285) & (0.234) & (0.254) & (0.247) &  & (0.239) & (3.102) & \\
Mean prior P(Black) &  &  &  & -2.128** &  & -2.248** & -1.266 & -2.202** & -2.179**\\
 &  &  &  & (0.840) &  & (0.847) & (0.910) & (0.839) & (0.852)\\
Residual roster share Black (net of position mix) &  &  &  &  &  & -0.470* &  &  & \\
 &  &  &  &  &  & (0.246) &  &  & \\
Position-mix expected share Black &  &  &  &  &  &  & -4.642** &  & \\
 &  &  &  &  &  &  & (1.987) &  & \\
Roster share Black, squared &  &  &  &  &  &  &  & 2.347 & \\
 &  &  &  &  &  &  &  & (2.380) & \\
Three-group Blau index (diversity) &  &  &  &  &  &  &  &  & 0.492*\\
 &  &  &  &  &  &  &  &  & (0.274)\\
Lagged win pct. &  &  & -0.078 & -0.145* & -0.089 & -0.145* & -0.146 & -0.149* & -0.143*\\
 &  &  & (0.094) & (0.085) & (0.093) & (0.084) & (0.088) & (0.087) & (0.084)\\
Lagged market expected wins &  &  & 0.031*** & 0.029*** & 0.030*** & 0.029*** & 0.029*** & 0.029*** & 0.029***\\
 &  &  & (0.009) & (0.010) & (0.009) & (0.010) & (0.010) & (0.010) & (0.010)\\
Opening-day coaches' share Black &  &  &  & -0.364 & -0.388* & -0.365 & -0.304 & -0.389 & -0.346\\
 &  &  &  & (0.241) & (0.215) & (0.240) & (0.223) & (0.246) & (0.240)\\
Opening-day head coach Black &  &  &  & -0.123 & -0.086 & -0.124 & -0.103 & -0.122 & -0.122\\
 &  &  &  & (0.089) & (0.080) & (0.088) & (0.094) & (0.090) & (0.087)\\
Opening-day starting QB Black &  &  &  & -0.041 & 0.012 & -0.047 & -0.031 & -0.044 & -0.047\\
 &  &  &  & (0.034) & (0.033) & (0.034) & (0.034) & (0.033) & (0.034)\\
\midrule
Observations & 384 & 384 & 384 & 384 & 384 & 384 & 384 & 384 & 384\\
R$^2$ & 0.038 & 0.273 & 0.353 & 0.421 & 0.367 & 0.418 & 0.434 & 0.423 & 0.417\\
Within R$^2$ & 0.038 & 0.025 & 0.132 & 0.223 & 0.151 & 0.219 & 0.241 & 0.226 & 0.219\\
WCB p-value, diversity & 0.002 & 0.049 & 0.015 & 0.032 & 0.008 & 0.066 & 0.007 & 0.345 & 0.093\\
Mean of outcome & 0.500 & 0.500 & 0.500 & 0.500 & 0.500 & 0.500 & 0.500 & 0.500 & 0.500\\
SD of diversity net of FE & 0.054 & 0.044 & 0.044 & 0.044 & 0.044 & 0.045 & 0.044 & 0.044 & 0.041\\
Effect of 1-SD change in expected share (win pct. pts.) & -3.71 & -2.63 & -2.75 & -2.48 & -3.04 & -2.11 & -3.03 & -2.42 & 2.02\\
Turning point (share) &  &  &  &  &  &  &  & 0.776 & \\
Season FE & Yes & Yes & Yes & Yes & Yes & Yes & Yes & Yes & Yes\\
Franchise FE & No & Yes & Yes & Yes & Yes & Yes & Yes & Yes & Yes\\
Lagged outcomes and roster quality & No & No & Yes & Yes & Yes & Yes & Yes & Yes & Yes\\
Calibration controls (priors, prior-covariate shares) & No & No & No & Yes & No & Yes & Yes & Yes & Yes\\
\bottomrule
\multicolumn{10}{l}{\rule{0pt}{1em}* p $<$ 0.1, ** p $<$ 0.05, *** p $<$ 0.01}\\
\end{tabular}}
\end{table}

{\footnotesize
\noindent\textit{Notes:} This table includes the estimation results of equation (2). The unit of observation is a franchise-season, 2014-2025 (384 franchise-seasons); the dependent variable is the regular-season win percentage (0-1). The snap-weighted share Black weights each player by his offensive plus defensive snaps for the franchise in the regular season; it is measured during season $t$ and can respond to injuries and benching. Snap weights also respond to game script: teams that trail run more offensive plays, teams with weak offenses play more defensive snaps, and defensive units have higher Black shares, so losing can raise the snap-weighted share mechanically (Table \ref{tab:roster-diversity-robustness} reports the unit-balanced share and a control for the defensive snap share). The coefficients on this share are therefore descriptive conditional associations between in-season playing time by race and results, not effects of a roster fixed before the season; the opening-day roster treatments of Table \ref{tab:roster-diversity-robustness}, columns (3)-(4), are the ones fixed before the season's results arrive. Columns with roster quality also include the prior season's cap share (the sum of players' cap percent on the franchise's season $t-1$ cap table; first observed in 2014, which is why the panel starts there). The season's own cap share is an annual total that in-season signings, releases and restructures change, so it is not fixed by opening day and enters only the contemporaneous column of Table \ref{tab:roster-diversity-robustness}. Column (1) includes season fixed effects; column (2) adds franchise fixed effects, so identification comes from changes in a franchise's roster composition over time. The identifying assumption is that, conditional on franchise and season fixed effects and the controls, within-franchise changes in roster composition are uncorrelated with other time-varying determinants of performance (talent, injuries, coaching). Column (3) adds lagged outcomes and opening-day roster quality, to address talent: rosters with different composition may simply be more talented. With franchise fixed effects over a short panel the lagged win percentage is subject to Nickell bias; it is a rough control for persistent talent, which is why lagged market expected wins enter as well. Column (4) adds the staff and quarterback controls, so the roster share is not proxying for a Black head coach, staff composition or the quarterback's race. Controls are fixed by opening day: the lagged win percentage and lagged market expected wins; roster quality of the active roster of the first regular-season game (mean log draft pick, undrafted = log 300; share of first-round picks; mean age; mean experience); and the Black share of the on-field coaches listed in the preseason staff snapshot, used only where the full staff is observed (2007 on; a missing indicator absorbs earlier seasons), the opening-day head coach's race and the opening-day starting quarterback's race. Under the predicted measures these race controls are predicted probabilities like the treatment, and the calibration controls include their priors (the coaches' mean prior and the head coach's and quarterback's priors). The head coach's predicted probability is calibrated to the staff population, not to the selected population of head coaches, so it is a noisy control. Under the predicted measure column (4), the headline specification, also includes the calibration controls: the mean prior P(Black), the weighted mean over players of the predicted-race model's prior probability of being Black, and the weighted shares of the levels of the prior's non-position covariates (draft round, college type, rookie era, hometown-county availability). The probabilities are posteriors given these covariates, so the coefficient on the expected share identifies the effect of the true share only conditional on them; draft round and college type also measure talent directly. Columns (6)-(9) include them as well. Column (5) drops them, as a sensitivity check. Net of franchise and season fixed effects the mean prior's correlation with the  share is 0.20. Placebo tests (Table \ref{tab:roster-diversity-placebo}): with the column (4) controls, next season's snap-weighted share enters the win percentage with -0.514 (SE 0.240; WCB p = 0.039) against -0.202 (SE 0.325; WCB p = 0.545) for the current share; for the headcount share the lead is -0.412 (SE 0.193; WCB p = 0.044) and the current share -0.178 (SE 0.237; WCB p = 0.459). Last season's win percentage predicts the current share with -0.0209 (SE 0.0096; WCB p = 0.035) (snap-weighted) and -0.0159 (SE 0.0079; WCB p = 0.055) (headcount); net of the current and lagged mean prior these become -0.0209 (SE 0.0097; WCB p = 0.041) and -0.0086 (SE 0.0066; WCB p = 0.208). Next season's composition is more strongly related to this season's results than the current composition, so the identifying assumption is rejected and these estimates should not be read as causal effects: performance shapes composition, or a persistent franchise trait drives both. Column (6) replaces the share with the residual share, to address positional composition: Black shares differ sharply by position (position at entry is also a covariate of the race prior). The residual share is the actual share minus the expected share given the team's position mix (the sum over position groups of the team's position weight times the league-season Black share at that position on all other franchises). Column (7) instead controls for the position-mix expected share; the leave-one-franchise-out league share makes it mechanically, but only slightly (of order 1/31 of the own-share variation), lower when the own share is high. The Black share measures composition, not diversity. Column (8) adds the squared share; it nests the two-group (Black/non-Black) Blau index $2s(1-s)$, which is a quadratic in the share. Its WCB p-value refers to the squared term; the effect row reports the marginal effect at the sample mean share, and the turning point is $-\beta_1/(2\beta_2)$. Column (9) uses the three-group (Black, white, other) Blau index $1 - b^2 - w^2 - o^2$, the probability that two players drawn at random belong to different groups, as the diversity measure. Net of franchise and season fixed effects its correlation with the Black share is -0.98; the mean share is 0.66, above one half, so a higher Black share means a less diverse roster. Under the predicted measures the index is computed from the expected shares (mean member probabilities), not as the expected index. The SD row is the SD of the diversity measure net of the column's fixed effects in the estimation sample; the effect row multiplies it by the coefficient and by 100 (win-percentage points). Roster-quality and calibration-control coefficients other than the mean prior are not shown. Standard errors, in parentheses, are clustered at the franchise level (32 clusters); wild cluster bootstrap-t p-values (restricted, WCR) use Webb weights and 9,999 replications. Under the model-only predicted measure a share Black is the expected share of players who are non-Hispanic Black alone (multiracial and Hispanic Black players count as non-Black); the documented, provisional and hand-coded sensitivity measures count Black alone or in combination. Under regression calibration the coefficient on an expected roster share (the mean over players of the predicted probability of being non-Hispanic Black) equals the effect of the true share only if the probabilities are calibrated at the team level given the regression's controls; this requires that players' names and hometown counties be unrelated to team performance given race and the prior's covariates, and the calibration controls (the team mean of the prior and the weighted shares of its draft-round, college-type, rookie-era and hometown-county levels, plus the priors of the coaches, head coach and quarterback whose predicted race enters as a control) condition on those covariates. Shrinkage toward the prior by itself causes no attenuation (the error of an expected share is of the Berkson type), but individual-level calibration does not imply team-level calibration: if the true share rises more (less) than one-for-one with the expected share across teams, the coefficient overstates (understates) the effect, so the sign of any bias is not known. A league-wide level error, such as the model's shortfall against the TIDES league share, is absorbed by the season fixed effects. A team-level calibration check among documented players (validation only): regressing a franchise-season's documented Black share of its documented players on the same players' mean predicted probability, with franchise and season fixed effects, gives a slope of 0.79 (SE 0.06) for snap weights and  0.69 (SE 0.06) for headcount weights, against one under calibration. Documented players are a selected, mostly Black and famous subset (about 28 percent of the weight), so the check does not establish the team-level calibration of the full roster in either direction. Effect rows refer to a one-SD change in the expected share, whose SD is smaller than that of the true share. Race is predicted, not observed: each person's probability of being non-Hispanic Black alone combines first name, surname and hometown county (BIFSG) with an NFL-specific prior estimated by EM on predetermined characteristics (players: position at entry, rookie era, draft round, college type, county availability; staff: role, unit and era at first appearance); documented race statements are not used (notes/race-prediction-design.md). Black Hispanic and multiracial Black persons are in the other-race probability, so P(Black) is narrower than the Black-alone-or-in-combination definition of the hand codes and of TIDES. Team shares are expected shares (mean member probability). Their error is Berkson-type: estimates are consistent if the probabilities are calibrated given the regression's controls (the team mean of the prior is one of them), but less precise than with observed race; individual coaches' probabilities are noisy because names carry little information for many Black coaches. The validation exhibit compares the predictions with published TIDES shares.
\par}

